\documentclass[aps,prd,twocolumn,superscriptaddress,nofootinbib]{revtex4-2}

\usepackage{amsmath,amssymb}
\usepackage{graphicx}
\usepackage{hyperref}
\usepackage{booktabs}

\begin{document}

\title{Link-graph spectral dimension on causal sets is determined by connectivity,\\not causal structure}

\author{Matt Loftus}
\email{matthew.a.loftus@gmail.com}
\affiliation{Independent researcher}

\date{\today}

\begin{abstract}
The spectral dimension $d_s$ flowing from 4 at large scales to $\sim$2 at the Planck scale is widely cited as a universal prediction of quantum gravity, appearing independently in causal dynamical triangulations (CDT), asymptotic safety, Ho\v{r}ava-Lifshitz gravity, loop quantum gravity, and causal set theory. We show that (i) the \emph{crossing} of $d_s = 2$ at some diffusion scale is a generic property of any finite graph, including random graphs with no gravitational content; (ii) the \emph{plateau} at $d_s \approx 2$ that distinguishes CDT from random graphs grows with system size in CDT but shrinks in causal sets; and (iii) the peak spectral dimension on causal set link graphs is purely determined by link density---random graphs with matched connectivity produce identical $d_s$ values. These results establish that link-graph spectral dimension is not a reliable probe of causal geometry in discrete quantum gravity. We propose plateau duration and structure-sensitive observables such as interval entropy as alternatives.
\end{abstract}

\maketitle

\section{Introduction}

The spectral dimension $d_s$, defined via the return probability of a diffusion process, has emerged as a key observable in quantum gravity. The seminal observation by Ambj\o{}rn, Jurkiewicz, and Loll~\cite{ambjorn2005} that four-dimensional causal dynamical triangulations (CDT) produce a spectral dimension flowing from $d_s = 4$ at large scales to $d_s \approx 1.8$ at short distances sparked enormous interest. Carlip~\cite{carlip2017} cataloged at least eight independent approaches that reproduce this dimensional reduction, calling it ``perhaps the most robust prediction of quantum gravity.''

In the causal set approach~\cite{bombelli1987,surya2019}, spectral dimension has been measured on Poisson sprinklings into known spacetimes~\cite{eichhorn2014}, but its behavior in dynamical simulations---such as those weighted by the Benincasa-Dowker action~\cite{benincasa2010,surya2012,glaser2018}---has not been systematically studied. Here we present such a study and find that link-graph spectral dimension on causal sets does not reliably probe causal geometry.

\section{Method}

The spectral dimension is defined via the normalized graph Laplacian. For an undirected graph with adjacency matrix $A$ and degree matrix $D$:
\begin{equation}
L = I - D^{-1/2} A \, D^{-1/2},
\end{equation}
with eigenvalues $\{\lambda_k\}$. The return probability is $P(\sigma) = n^{-1} \sum_k e^{-\lambda_k \sigma}$, and the spectral dimension is
\begin{equation}
d_s(\sigma) = -2 \frac{d \ln P}{d \ln \sigma}.
\end{equation}
For causal sets, we symmetrize the link matrix to obtain the undirected adjacency. We use this identical procedure for all graph types, enabling direct comparison.

We compare six graph families at matched system sizes ($N = 100$--$800$): (i) causal sets sprinkled into 2D Minkowski, (ii) 2D CDT configurations sampled by MCMC, (iii) classical sequential growth (CSG) causal sets, (iv) random directed acyclic graphs (DAGs), (v) Erd\H{o}s-R\'enyi random graphs, and (vi) regular lattices.

\section{The crossing theorem}

The $d_s = 2$ crossing is not an empirical observation---it is a mathematical inevitability.

\textbf{Proposition.} \emph{For any finite connected graph, $d_s(\sigma)$ is smooth on $(0,\infty)$ with $d_s \to 0$ as $\sigma \to 0^+$ and $d_s \to 0$ as $\sigma \to \infty$. If $d_s^{\max} > c$ for any $c > 0$, the intermediate value theorem gives $\sigma_1 < \sigma_2$ with $d_s(\sigma_1) = d_s(\sigma_2) = c$.}

\emph{Proof.} For a connected graph with eigenvalues $0 = \lambda_0 < \lambda_1 \leq \cdots \leq \lambda_{n-1}$:
\begin{equation}
d_s(\sigma) = \frac{2\sigma}{nP(\sigma)} \sum_{k=1}^{n-1} \lambda_k e^{-\lambda_k\sigma}.
\end{equation}
Smoothness follows from $P > 0$ (the $k = 0$ term contributes $1/n$). As $\sigma \to 0^+$: $P \to 1$, $\sigma P' \to 0$, so $d_s \to 0$. As $\sigma \to \infty$: $P \to 1/n$, $|\sigma P'| \leq \sigma\lambda_{\max}(n{-}1)e^{-\lambda_1\sigma} \to 0$ (since $\lambda_1 > 0$), so $d_s \to 0$. The IVT gives the result. $\square$

\textbf{Remark.} The value $c = 2$ plays no special role. The crossing is a topological consequence of the IVT, not a signature of dimensional reduction.

At fixed $n$, the peak $d_s^{\max}$ increases monotonically with mean degree $\bar{k}$, because higher connectivity concentrates the Laplacian spectrum (bandwidth $\Delta\lambda \sim 4/\sqrt{\bar{k}}$ for Erd\H{o}s-R\'enyi graphs), producing a sharper return-probability peak. This explains why random graphs with matched link density yield the same $d_s$ as causal sets.

\section{Numerical results}

\subsection{Verification of the crossing theorem}

Every graph type tested crosses $d_s = 2$ (Fig.~\ref{fig:peaks}, Table~\ref{tab:crossing}), as guaranteed by the proposition.

\begin{figure}[t]
\includegraphics[width=\columnwidth]{fig1_peaks.pdf}
\caption{Peak spectral dimension for different graph types at $N = 500$. The dashed line marks $d_s = 2$. Random graphs (red) produce peak values comparable to causal sets (blue/green). Only lattices (orange) correctly reflect their true dimension.}
\label{fig:peaks}
\end{figure}

\begin{table}[h]
\centering
\begin{tabular}{lcc}
\toprule
Graph type & Peak $d_s$ & Physical? \\
\midrule
CSG causet ($p = 0.05$) & 4.81 & QG dynamics \\
Sprinkled 2D causet & 5.21 & Known manifold \\
Sprinkled 4D causet & 6.61 & Known manifold \\
Random DAG ($p = 0.01$) & 4.93 & None \\
Erd\H{o}s-R\'enyi ($p = 0.01$) & 4.95 & None \\
1D lattice & 1.21 & $d = 1$ \\
2D lattice ($22 \times 22$) & 2.29 & $d = 2$ \\
\bottomrule
\end{tabular}
\caption{Peak spectral dimension for different graph types at $N = 500$. Random graphs with no gravitational content produce peak $d_s$ comparable to causal sets. Only the lattices correctly reflect their true dimension.}
\label{tab:crossing}
\end{table}

\subsection{CDT plateau vs.\ causet crossing}

The meaningful signal is not the crossing but the \emph{plateau}: a range of scales over which $d_s$ remains near 2. We define the plateau width as the fraction of sampled $\sigma$ values for which $|d_s - 2| < 0.3$. (Results are qualitatively unchanged for thresholds 0.2 and 0.5; the CDT plateau grows and the causet plateau shrinks regardless of the threshold choice.)

\begin{table}[h]
\centering
\begin{tabular}{lccccc}
\toprule
& \multicolumn{2}{c}{CDT} & & \multicolumn{2}{c}{Causet 2D} \\
\cmidrule{2-3} \cmidrule{5-6}
$N$ & Peak $d_s$ & Plateau & & Peak $d_s$ & Plateau \\
\midrule
100 & 2.49 & 17\% & & 3.64 & 8\% \\
200 & 2.45 & 30\% & & 4.30 & 7\% \\
400 & 2.50 & 36\% & & 4.98 & 7\% \\
800 & 2.59 & 47\% & & 5.66 & 6\% \\
\bottomrule
\end{tabular}
\caption{Finite-size scaling of the spectral dimension plateau. CDT's plateau \emph{grows} with $N$ (17\%$\to$47\%), while the causet's ``plateau'' \emph{shrinks} (8\%$\to$6\%). CDT's peak $d_s$ is stable near 2.5 (correctly 2D); the causet peak diverges.}
\label{tab:plateau}
\end{table}

\begin{figure}[t]
\includegraphics[width=\columnwidth]{fig2_scaling.pdf}
\caption{Finite-size scaling of the spectral dimension plateau width (left) and peak value (right). CDT's plateau grows with $N$ while the causet ``plateau'' shrinks. CDT's peak is stable near the true $d = 2$; the causet peak diverges.}
\label{fig:scaling}
\end{figure}

Table~\ref{tab:plateau} and Fig.~\ref{fig:scaling} reveal a stark contrast. In CDT, the plateau width grows monotonically with system size, from 17\% at $N = 100$ to 47\% at $N = 800$, and the peak $d_s$ remains stable near $2.5$---correctly reflecting the 2D target. In sprinkled 2D causal sets, the plateau width \emph{decreases} from 8\% to 6\%, and the peak $d_s$ \emph{diverges} from the true dimension, reaching 5.66 at $N = 800$.

This demonstrates that $d_s \approx 2$ in CDT is a genuine, growing plateau (physical signal), while in causal sets it is a transient crossing that shrinks with system size (artifact).

\subsection{Spectral dimension is determined by link density}

To test whether $d_s$ on causal sets captures causal structure or merely connectivity, we compare causal sets sampled from the Benincasa-Dowker partition function against random graphs with matched link density.

In the continuum phase of the BD transition (link density $L/N \approx 2.6$), the peak spectral dimension is $d_s \approx 3.0$. In the crystalline phase ($L/N \approx 10.6$), $d_s \approx 4.1$. However, random graphs with the same link densities yield essentially identical values: $d_s \approx 3.0$ at $L/N = 2.6$ and $d_s \approx 4.0$ at $L/N = 10.6$ (Table~\ref{tab:control}).

\begin{table}[h]
\centering
\begin{tabular}{lcc}
\toprule
Link density $L/N$ & BD causet $d_s$ & Random graph $d_s$ \\
\midrule
2.6 (continuum) & 3.0 & 3.0 \\
5.0 & --- & 3.6 \\
8.0 & --- & 3.9 \\
10.6 (crystalline) & 4.1 & 4.0 \\
\bottomrule
\end{tabular}
\caption{Peak spectral dimension vs.\ link density. Random graphs with matched $L/N$ produce identical $d_s$, demonstrating that link-graph spectral dimension is purely a function of connectivity.}
\label{tab:control}
\end{table}

This result is definitive: the spectral dimension computed from the link graph of a causal set is determined by the number of links, not by causal structure. It cannot distinguish the continuum phase from the crystalline phase of the BD transition, nor can it distinguish a causal set from a random graph with the same connectivity.

\section{Discussion}

Our findings have several implications for the quantum gravity program:

\textbf{1. Claims of $d_s \to 2$ require qualification.} The crossing of $d_s = 2$ is inevitable for any finite graph with peak $d_s > 2$. Studies claiming dimensional reduction to 2 in a given approach should report the plateau duration (scale range over which $|d_s - 2| < \epsilon$) and its scaling with system size, not merely the existence of a scale where $d_s = 2$.

\textbf{2. CDT's dimensional reduction is genuine.} The CDT plateau grows with system size, which is the expected behavior for a physical signal that sharpens in the continuum limit. This is consistent with Carlip's argument~\cite{carlip2017} that $d_s = 2$ reflects the dimensionless nature of Newton's constant in 2D, connected to BKL-type ``asymptotic silence.'' We emphasize that our comparison is in 2D; the famous $d_s$ flow from 4 to $\sim$2 occurs in 4D CDT, and extending our analysis to higher dimensions is an important direction for future work.

\textbf{3. The problem is not the operator.} One might hope that replacing the link-graph Laplacian with the Benincasa-Dowker d'Alembertian~\cite{benincasa2010}---which is specifically designed to approximate the continuum wave operator---would fix the spectral dimension. We tested this and found that the BD d'Alembertian also fails the random graph control: sprinkled 2D causets and random DAGs with matched connectivity produce identical $d_s$ values ($3.2 \pm 0.1$ vs.\ $3.3 \pm 0.1$). The problem is not the choice of diffusion operator but the discrete causal structure itself, which lacks the regularity needed for diffusion-based dimension estimation. The Myrheim-Meyer estimator~\cite{myrheim1978,meyer1988} and interval entropy (the Shannon entropy of the order-interval size distribution) are robust alternatives that capture causal structure independently of link density.

\textbf{4. Reinterpretation of the Eichhorn-Mizera anomaly.} Eichhorn and Mizera~\cite{eichhorn2014} found that spectral dimension on causal set link graphs \emph{increases} at small scales---the opposite of CDT. They attributed this to the ``non-locality'' of causal sets arising from Lorentz-invariant sprinkling, which produces unbounded link valency (each element has infinitely many nearest neighbors in the continuum limit). Our link-density finding provides a simpler explanation: at small diffusion scales, the walker probes the high-connectivity local link neighborhood, and $d_s$ faithfully tracks this connectivity rather than manifold dimension.

This reinterpretation is consistent with the subsequent literature. Belenchia \emph{et al.}~\cite{belenchia2016} recovered $d_s \to 2$ using the continuum nonlocal d'Alembertian (bypassing the link graph entirely). Eichhorn, Surya, and Versteegen~\cite{eichhorn2019} recovered dimensional reduction by restricting to spatial hypersurfaces, which bounds the link valency. In our framework, all these fixes work for the same reason: they tame or bypass the scale-dependent connectivity of the link graph.

\textbf{5. Random graph controls are essential.} Any observable proposed as a probe of emergent geometry should be tested against random graphs with matched connectivity. If the observable cannot distinguish structured from random graphs at the same link density, it does not capture geometry.

\section{Conclusion}

Link-graph spectral dimension on causal sets is determined by connectivity, not causal structure. The $d_s \approx 2$ crossing is generic to all finite graphs, while the $d_s \approx 2$ plateau is specific to CDT (where it grows with system size). We recommend that future studies of spectral dimension in discrete quantum gravity (i) report plateau duration and its scaling, (ii) test against random graph controls, and (iii) consider structure-sensitive alternatives such as interval entropy for causal sets.

\begin{acknowledgments}
This work was conducted as an independent computational exploration with assistance from Claude (Anthropic).
\end{acknowledgments}

\begin{thebibliography}{99}

\bibitem{ambjorn2005}
J.~Ambj\o{}rn, J.~Jurkiewicz, and R.~Loll,
``Spectral dimension of the universe,''
Phys.\ Rev.\ Lett.\ \textbf{95}, 171301 (2005),
\href{https://arxiv.org/abs/hep-th/0505113}{arXiv:hep-th/0505113}.

\bibitem{carlip2017}
S.~Carlip,
``Dimension and dimensional reduction in quantum gravity,''
Class.\ Quant.\ Grav.\ \textbf{34}, 193001 (2017),
\href{https://arxiv.org/abs/1705.05417}{arXiv:1705.05417}.

\bibitem{bombelli1987}
L.~Bombelli, J.~Lee, D.~Meyer, and R.~D.~Sorkin,
``Space-time as a causal set,''
Phys.\ Rev.\ Lett.\ \textbf{59}, 521 (1987).

\bibitem{surya2019}
S.~Surya,
``The causal set approach to quantum gravity,''
Living Rev.\ Rel.\ \textbf{22}, 5 (2019),
\href{https://arxiv.org/abs/1903.11544}{arXiv:1903.11544}.

\bibitem{eichhorn2014}
A.~Eichhorn and S.~Mizera,
``Spectral dimension in causal set quantum gravity,''
Class.\ Quant.\ Grav.\ \textbf{31}, 125007 (2014),
\href{https://arxiv.org/abs/1311.2530}{arXiv:1311.2530}.

\bibitem{benincasa2010}
D.~M.~T.~Benincasa and F.~Dowker,
``The scalar curvature of a causal set,''
Phys.\ Rev.\ Lett.\ \textbf{104}, 181301 (2010),
\href{https://arxiv.org/abs/1001.2725}{arXiv:1001.2725}.

\bibitem{surya2012}
S.~Surya,
``Evidence for the continuum in 2D causal set quantum gravity,''
Class.\ Quant.\ Grav.\ \textbf{29}, 132001 (2012),
\href{https://arxiv.org/abs/1110.6244}{arXiv:1110.6244}.

\bibitem{glaser2018}
L.~Glaser, D.~O'Connor, and S.~Surya,
``Finite size scaling in 2d causal set quantum gravity,''
Class.\ Quant.\ Grav.\ \textbf{35}, 045006 (2018),
\href{https://arxiv.org/abs/1706.06432}{arXiv:1706.06432}.

\bibitem{myrheim1978}
J.~Myrheim,
``Statistical geometry,''
CERN preprint TH-2538 (1978).

\bibitem{meyer1988}
D.~A.~Meyer,
``The dimension of causal sets,''
Ph.D. thesis, MIT (1988).

\bibitem{belenchia2016}
A.~Belenchia, D.~M.~T.~Benincasa, A.~Marcian\`o, and L.~Modesto,
``Spectral dimension from nonlocal dynamics on causal sets,''
Phys.\ Rev.\ D \textbf{93}, 044017 (2016),
\href{https://arxiv.org/abs/1507.00330}{arXiv:1507.00330}.

\bibitem{eichhorn2019}
A.~Eichhorn, S.~Surya, and F.~Versteegen,
``Spectral dimension on spatial hypersurfaces in causal set quantum gravity,''
Class.\ Quant.\ Grav.\ \textbf{36}, 205013 (2019),
\href{https://arxiv.org/abs/1905.13498}{arXiv:1905.13498}.

\end{thebibliography}

\end{document}
